\section{Unsupervised learning}
\label{sec:unsupervised}

Everything so far has been supervised: every input $x_i$ came with a label $y_i$, and
the whole machinery was built to predict $y$ from $x$. We now remove the labels.

\subsection{The setup, and what changes without labels}
\label{subsec:unsup-setup}

The data is a plain collection of points
\[
x_1, \dots, x_n \in \R^d ,
\]
with nothing attached to them, and we stack them as always into the design matrix
$X \in \R^{n \times d}$ whose $i$-th row is $x_i$. There is no $y$, and therefore
no function $f$ to fit and no prediction to be scored. The question becomes the vaguer
one of \emph{finding structure}: is the cloud of points really spread over all $d$
directions, or does it essentially live in a much smaller subspace? Does it break up
into a few well-separated groups?

\paragraph{What we keep: one recipe, two instances.}
The good news is that the main ideas of these notes stay mostly intact. We still
write down a class of simplified descriptions of the data, still measure quality by a
squared error, and still minimise. The unifying statement, which is the thread running
through this whole section, is that both methods we study solve
\[
\ \min_{\hat X \,\in\, \mathcal{C}} \ \big\Vert X - \hat X \big\Vert_F^2 
\]
for two different choices of the set $\mathcal{C}$ of allowed reconstructions
$\hat X \in \R^{n \times d}$. In both cases $\hat X$ is a \emph{simplified stand-in}
for the data: row $i$ of $\hat X$ is what we keep of the point $x_i$, and the objective
asks that we lose as little as possible.
\begin{itemize}
    \item \emph{Principal component analysis} takes $\mathcal{C}$ to be the matrices
    whose rows all lie in one $K$-dimensional linear subspace of $\R^d$. Each point is
    replaced by its projection onto a common subspace: the data is compressed in
    \emph{direction}.
    \item \emph{$K$-means} takes $\mathcal{C}$ to be the matrices whose rows take at
    most $K$ distinct values. Each point is replaced by the centroid of its group: the
    data is compressed in \emph{identity}, since after the reconstruction all we
    remember of $x_i$ is which of $K$ representatives it was closest to.
\end{itemize}
The parameter $K$ measures, in both cases, how much we allow ourselves to keep, and it
is a hyperparameter we must pick: it is not fitted by the
optimiser.

\begin{remark}[The Frobenius norm]
\label{rem:frobenius}
The norm above is the \emph{Frobenius} norm, which is nothing more exotic than the
Euclidean norm of a matrix read as a long vector:
\[
\Vert A \Vert_F^2 \ \coloneqq \ \sum_{i,j} A_{i,j}^2 \ = \ \tr(A^\top A) .
\]
When the rows of $A$ are $a_1, \dots, a_n$ this is $\sum_{i=1}^n \Vert a_i \Vert^2$, so
$\Vert X - \hat X \Vert_F^2 = \sum_{i=1}^n \Vert x_i - \hat x_i \Vert^2$ really is the
total squared reconstruction error, summed over the samples. We use throughout the
\emph{cyclicity of the trace}, $\tr(AB) = \tr(BA)$ whenever both products are defined;
see \Cref{sec:cheat-linalg}.
\end{remark}

\begin{remark}[A third index]
\label{rem:third-index}
Both methods need to count \emph{groups} (clusters, or directions kept), a third kind
of index alongside the samples $i$ and the coordinates $j$ used throughout these notes. We
use $k$ for it, ranging over $1, \dots, K$, exactly as $k$ already ranged over the
classes in \Cref{sec:classification}. So $x_i$ is a sample, $x_{i,j}$ one of its
coordinates, and $\mu_k$ or $v_k$ the $k$-th centroid or direction.
\end{remark}


\subsection{Principal component analysis}
\label{subsec:pca}

\paragraph{Centring.}
PCA is about \emph{directions of variation}, so the position of the cloud is
irrelevant and must be removed first. Throughout this subsection we assume the data has
been \emph{centred}: writing $\bar x = \tfrac1n \sum_{i=1}^n x_i$ for the empirical
mean, we replace each $x_i$ by $x_i - \bar x$, so that
\[
\frac{1}{n} \sum_{i=1}^n x_i = 0 ,
\]
and we keep calling the resulting matrix $X$. This is an important step: without it
the first principal direction is dominated by the direction of $\bar x$, which says
where the data sits rather than how it varies. Whether one should further
\emph{standardise} the coordinates is a separate and more debatable question, discussed
at the end of the subsection.

\paragraph{Projecting onto a subspace.}
We look for a linear subspace $S \subseteq \R^d$ of dimension $K \leq d$ on which to
project the data. Describe it by an orthonormal basis $v_1, \dots, v_K$, that is,
$\langle v_k, v_{k'} \rangle = \delta_{k,k'}$, and gather the basis into a matrix
\[
V = \big( v_1 \mid \dots \mid v_K \big) \in \R^{d \times K} ,
\qquad \text{so that} \qquad
V^\top V = I_K ,
\]
the last identity being exactly the statement that the columns are orthonormal.
Two matrices are then worth reading carefully, because everything follows from them.

The matrix $XV \in \R^{n \times K}$ collects the \emph{coordinates of the data in the
new basis}:
\[
XV = \begin{pmatrix}
\langle x_1, v_1 \rangle & \dots & \langle x_1, v_K \rangle \\
\vdots & & \vdots \\
\langle x_n, v_1 \rangle & \dots & \langle x_n, v_K \rangle
\end{pmatrix} .
\]
Its $i$-th row is the compressed version of $x_i$: $K$ numbers instead of $d$. These
are the \emph{scores}, or \emph{principal components}, of the data.

The matrix $X V V^\top \in \R^{n \times d}$ puts them back into the original space:
\[
X V V^\top = \begin{pmatrix}
- & \sum_{k=1}^K \langle x_1, v_k \rangle\, v_k & - \\
  & \vdots & \\
- & \sum_{k=1}^K \langle x_n, v_k \rangle\, v_k & -
\end{pmatrix} ,
\]
whose $i$-th row is the orthogonal projection of $x_i$ onto $S$. So $P_V \coloneqq
V V^\top \in \R^{d \times d}$ is the orthogonal projector onto $S$, and $\hat X = X
V V^\top$ is precisely the reconstruction of \Cref{subsec:unsup-setup}, with
$\mathcal{C}$ the set of matrices whose rows lie in a common $K$-dimensional subspace.

\begin{exercise}[$\star$]
Check that $P_V = VV^\top$ satisfies $P_V^\top = P_V$ and $P_V^2 = P_V$, so it is
indeed an orthogonal projector, and that $\rank(P_V) = K$. Where is $V^\top V = I_K$
used? Note that $V V^\top \neq I_d$ unless $K = d$: the order of the product matters
enormously.
\end{exercise}

\paragraph{The objective, and its two readings.}
PCA chooses the subspace minimising the reconstruction error:
\[
\min_{V \in \R^{d \times K}, \ V^\top V = I_K} \ \big\Vert X - X V V^\top \big\Vert_F^2 .
\]
The following computation, which is the heart of the matter, turns this into an
eigenvalue problem.

\begin{prop}[Reconstruction error and retained variance]
\label{prop:pca-two-readings}
For every $V \in \R^{d \times K}$ with $V^\top V = I_K$,
\[
\big\Vert X - X V V^\top \big\Vert_F^2
\ = \ \Vert X \Vert_F^2 \ - \ \tr\big( V^\top X^\top X\, V \big) .
\]
Since $\Vert X \Vert_F^2$ does not depend on $V$, \emph{minimising the reconstruction
error is the same problem as maximising} $\tr(V^\top X^\top X V)$.
\end{prop}

\begin{proof}
Expand the square using $\Vert A \Vert_F^2 = \tr(A^\top A)$:
\[
\big\Vert X - X V V^\top \big\Vert_F^2
= \tr(X^\top X) - 2\, \tr\big( X^\top X V V^\top \big)
+ \tr\big( V V^\top X^\top X V V^\top \big) ,
\]
where the cross terms were merged using $\tr(A^\top) = \tr(A)$. For the last term,
cyclicity of the trace moves the leading $V$ to the end,
\[
\tr\big( V V^\top X^\top X V V^\top \big) = \tr\big( V^\top X^\top X V \, V^\top V \big)
= \tr\big( V^\top X^\top X V \big) ,
\]
using $V^\top V = I_K$; and the same cyclicity gives $\tr(X^\top X V V^\top) =
\tr(V^\top X^\top X V)$. The two therefore combine into $-2 + 1 = -1$ times
$\tr(V^\top X^\top X V)$, which is the claim.
\end{proof}

The quantity being maximised has a direct statistical meaning. Since the data is
centred,
\[
\tr\big( V^\top X^\top X V \big) = \sum_{k=1}^K v_k^\top X^\top X\, v_k
= \sum_{k=1}^K \sum_{i=1}^n \langle x_i, v_k \rangle^2
= \big\Vert X V \big\Vert_F^2 ,
\]
which is $n$ times the total \emph{empirical variance of the projected data}. So the
two natural things one might have asked for turn out to be the same request:

\begin{center}
\begin{tabular}{@{}c@{\quad}c@{\quad}c@{}}
\emph{keep the subspace} & $\Longleftrightarrow$ & \emph{keep the subspace} \\
\emph{losing the least information} & & \emph{carrying the most variance.}
\end{tabular}
\end{center}

\noindent
This is the first thing to remember about PCA. It is also why the empirical covariance
matrix
\[
\hat \Sigma \ \coloneqq \ \frac{1}{n} X^\top X \ \in \R^{d \times d}
\]
is the object that decides everything: it is symmetric positive semi-definite (see
\Cref{sec:cheat-convexity}), and the problem is now to maximise $\tr(V^\top \hat\Sigma
V)$ over orthonormal $V$.

\paragraph{Solving it.}
Start with a single direction, where the answer is the classical Rayleigh quotient
bound.

\begin{lemma}[One direction]
\label{lem:rayleigh}
Let $A \in \R^{d \times d}$ be symmetric with eigenvalues $\lambda_1 \geq \dots \geq
\lambda_d$ and corresponding orthonormal eigenvectors $u_1, \dots, u_d$. Then
\[
\max_{\Vert v \Vert = 1} \ v^\top A v = \lambda_1 ,
\]
attained at $v = u_1$.
\end{lemma}

\begin{proof}
Write $v = \sum_{j=1}^d c_j u_j$ in the eigenbasis, so that $\Vert v \Vert^2 = \sum_j
c_j^2 = 1$. Then $v^\top A v = \sum_j \lambda_j c_j^2 \leq \lambda_1 \sum_j c_j^2 =
\lambda_1$, a weighted average of the eigenvalues being at most the largest. Equality
requires all the weight to sit on eigenvalues equal to $\lambda_1$, and $v = u_1$ works.
\end{proof}

The general case is the same statement one dimension at a time, and the proof is
pleasantly short once one thinks in terms of the projector rather than the basis.

\begin{thm}[PCA]
\label{thm:pca}
With $A = X^\top X$ symmetric positive semi-definite, with eigenvalues $\lambda_1 \geq
\dots \geq \lambda_d \geq 0$ and orthonormal eigenvectors $u_1, \dots, u_d$,
\[
\max_{V^\top V = I_K} \ \tr\big( V^\top X^\top X V \big) = \lambda_1 + \dots + \lambda_K ,
\]
attained by $V = (u_1 \mid \dots \mid u_K)$. The optimal subspace is therefore the span
of the $K$ leading eigenvectors of $X^\top X$, equivalently of the empirical covariance
$\hat\Sigma$, and the minimal reconstruction error is $\lambda_{K+1} + \dots +
\lambda_d$. The vectors $u_1, u_2, \dots$ are called the \emph{principal directions},
and $u_k$ the $k$-th one.
\end{thm}

\begin{proof}
Put $B = V V^\top$, the orthogonal projector onto the chosen subspace, and set
$\alpha_j \coloneqq u_j^\top B\, u_j$ for each $j$. Two observations:
\[
\alpha_j = \Vert V^\top u_j \Vert^2 \geq 0 ,
\qquad
\alpha_j \leq \Vert u_j \Vert^2 = 1 ,
\]
the second because an orthogonal projection never lengthens a vector, and
\[
\sum_{j=1}^d \alpha_j = \tr(B) = \tr(V V^\top) = \tr(V^\top V) = \tr(I_K) = K .
\]
Now expand $A = \sum_j \lambda_j u_j u_j^\top$ and use cyclicity again:
\[
\tr\big( V^\top A V \big) = \tr(A B) = \sum_{j=1}^d \lambda_j\, u_j^\top B u_j
= \sum_{j=1}^d \lambda_j \alpha_j .
\]
So the problem is bounded by the largest value of $\sum_j \lambda_j \alpha_j$ over all
$\alpha \in [0,1]^d$ with $\sum_j \alpha_j = K$. Since the $\lambda_j$ are sorted, the
best possible allocation of a total budget $K$ into weights capped at $1$ is to give a
full unit to each of the $K$ largest, giving $\lambda_1 + \dots + \lambda_K$. Taking
$V = (u_1 \mid \dots \mid u_K)$ gives exactly $\alpha_j = \mathbf{1}\{ j \leq K\}$, so
the bound is attained. Finally the minimal reconstruction error follows from
\Cref{prop:pca-two-readings} and $\Vert X \Vert_F^2 = \tr(X^\top X) = \sum_j \lambda_j$.
\end{proof}

\begin{remark}[The solution is nested, and the subspace is what matters]
The proof shows something the statement hides: the optimal $K$-dimensional subspace
\emph{contains} the optimal $(K-1)$-dimensional one. One does not re-solve the problem
for each $K$; one computes the eigenvectors once, sorted, and takes as many as one
wants. Note also that $V$ itself is not unique (any rotation of an orthonormal basis of
the same subspace gives the same reconstruction, and eigenvectors are only defined up
to sign), but the \emph{subspace} is, as soon as $\lambda_K > \lambda_{K+1}$.
\end{remark}

\begin{remark}[In practice one uses the SVD]
Forming $X^\top X$ and diagonalising it is the textbook route but not the numerical one,
for the reason given in \Cref{subsec:numerics}: it squares the condition number. Every
library instead computes the singular value decomposition $X = U D W^\top$, with $U \in
\R^{n \times n}$ and $W \in \R^{d \times d}$ orthogonal and $D \in \R^{n \times d}$ zero
off its diagonal, with diagonal entries $s_1 \geq s_2 \geq \dots \geq 0$. Then $X^\top X = W D^\top D W^\top$, so the principal
directions are the columns of $W$ (the right singular vectors) and $\lambda_k = s_k^2$.
The scores are $XV = U D$ restricted to the first $K$ columns.
\end{remark}

\paragraph{Choosing $K$.}
The eigenvalues answer this too, at least descriptively. The quantity
\[
\frac{\lambda_1 + \dots + \lambda_K}{\lambda_1 + \dots + \lambda_d} \ \in [0, 1]
\]
is the \emph{fraction of variance explained} by the first $K$ directions, and plotting
either it or the individual $\lambda_k$ against $k$ (a \emph{scree plot}) is the usual
way to decide. One looks for an elbow, a $K$ after which the eigenvalues fall off a
cliff, or simply takes the smallest $K$ explaining, say, $95\%$ of the variance. Both
are rules of thumb rather than theorems: there is no test error here to minimise. When
PCA is used as a preprocessing step before a supervised model, however, $K$ becomes an
ordinary hyperparameter of that model and can be chosen honestly by cross-validation.

\begin{remark}[Two warnings]
\emph{Centring is not optional, scaling is a modelling
choice.} PCA maximises variance, and variance has units: a feature measured in
millimetres has a variance $10^6$ times larger than the same feature in metres, and
will monopolise the first direction for no better reason. When the coordinates are
heterogeneous quantities one therefore standardises them first
(\Cref{subsec:standardise}), which amounts to diagonalising the empirical
\emph{correlation} matrix rather than the covariance. When all coordinates are
commensurate, as with the pixels of an image, standardising is usually the wrong thing
to do, since the differences in variance are then real information.
\end{remark}

\begin{exercise}[$\star$]
Take $d = 2$ and a cloud of points lying almost exactly on a line through the origin.
What are $\lambda_1, \lambda_2$ roughly, what is $u_1$, and what does the reconstruction
$X V V^\top$ look like for $K = 1$? Draw the picture.
\end{exercise}

\begin{exercise}[$\star\star$]
Show that PCA also solves the following apparently different problem: among all
matrices $\hat X$ of rank at most $K$, the one minimising $\Vert X - \hat X\Vert_F$ is
$\hat X = X V V^\top$ with $V$ as in \Cref{thm:pca}. (Start by checking that the rows of
$X V V^\top$ lie in a $K$-dimensional subspace, so $\rank(XVV^\top) \le K$; then argue
that any $\hat X$ of rank $\leq K$ has its rows in some $K$-dimensional subspace, and
that for a \emph{fixed} such subspace the best $\hat X$ is the orthogonal projection of
$X$ onto it.) This is a special case of the Eckart-Young theorem.
\end{exercise}


\subsection{$K$-means}
\label{subsec:kmeans}

We now compress in the other direction. Instead of asking that all the points be well
described by a common subspace, we ask that each point be well described by \emph{one
of $K$ representatives}.

\paragraph{The objective.}
Fix the hyperparameter $K$, the number of clusters. We look for
\begin{itemize}
    \item a \emph{partition} of the samples into $K$ groups $S_1, \dots, S_K$, i.e.\
    disjoint sets with $S_1 \cup \dots \cup S_K = \{1, \dots, n\}$;
    \item and $K$ \emph{centroids} $\mu_1, \dots, \mu_K \in \R^d$, one per group,
\end{itemize}
minimising the total squared distance from each point to the centroid of its own group,
\[
\mathcal{L}(S, \mu) \ = \ \sum_{k=1}^K \ \sum_{i \in S_k} \big\Vert x_i - \mu_k \big\Vert^2 .
\]
This is the natural way to say ``the groups should be tight''. Note that only the
squared Euclidean distance appears, which is already a modelling commitment: it is what
will make the clusters come out round.

\paragraph{Matrix form: the same problem as PCA, differently constrained.}
The partition is more conveniently encoded by an \emph{assignment matrix}. Let
\[
Z = (z_{i,k}) \in \{0, 1\}^{n \times K} , \qquad
\sum_{k=1}^K z_{i,k} = 1 \quad \text{for every } i ,
\]
so that $z_{i,k} = 1$ exactly when sample $i$ belongs to cluster $k$, each row of $Z$
containing a single $1$. Stack the centroids into
\[
M = \begin{pmatrix} - & \mu_1 & - \\ & \vdots & \\ - & \mu_K & - \end{pmatrix}
\in \R^{K \times d} .
\]
Then row $i$ of the product $ZM$ is $\sum_k z_{i,k} \mu_k$, which is simply the centroid
assigned to sample $i$, and therefore
\[
\mathcal{L}(M, Z)
= \sum_{i=1}^n \sum_{k=1}^K z_{i,k} \big\Vert x_i - \mu_k \big\Vert^2
= \big\Vert X - Z M \big\Vert_F^2 .
\]
So $K$-means is
\[
\min_{M \in \R^{K \times d}} \ \min_{\substack{Z \in \{0,1\}^{n \times K} \\ \sum_k z_{i,k} = 1 \ \forall i}}
\ \big\Vert X - Z M \big\Vert_F^2 ,
\]
which is exactly the reconstruction problem of \Cref{subsec:unsup-setup}: the
reconstruction is $\hat X = ZM$, whose rows take at most $K$ distinct values, namely the
$\mu_k$. Compare with PCA, where $\hat X = X V V^\top$ had rows in a $K$-dimensional
subspace. Same objective, same $K$, different notion of ``simple''.

\begin{exercise}[$\star$]
Check the middle equality above in detail: why does $\sum_k z_{i,k}\Vert x_i -
\mu_k\Vert^2$ equal $\Vert x_i - \sum_k z_{i,k}\mu_k\Vert^2$? (Use that exactly one term
of the sum is nonzero. Warning: this would be false for a general $Z$ with rows summing
to $1$ but entries in $[0,1]$ rather than $\{0,1\}$.)
\end{exercise}

\paragraph{The bad news.}
Unlike every problem we have met so far, this one is genuinely intractable: minimising
$\mathcal{L}$ jointly over $Z$ and $M$ is \textbf{NP-hard}, already for $d = 2$. The
reason is visible in the formulation: $Z$ ranges over a \emph{discrete} set, of size
$K^n$, and no amount of convexity is going to help, since the objective is not even
defined on a continuum. Exhaustive search is hopeless, and there is no closed form and
no gradient descent to fall back on.

\paragraph{The good news: alternating minimisation.}
What saves the situation is that although the joint problem is hard, each of the two
variables is easy \emph{given the other}, and in both cases the partial minimisation has
a closed form. That suggests the obvious algorithm: minimise alternately in $Z$ and in
$M$. This is \emph{Lloyd's algorithm}, universally called ``the'' $K$-means algorithm.

\begin{lemma}[The two steps]
\label{lem:kmeans-steps}
\emph{(Assignment.)} For a fixed $M$, the minimiser over admissible $Z$ is obtained
independently for each sample by assigning it to its nearest centroid:
\[
z_{i,k} = 1 \quad \text{if} \quad k = \argmin_{1 \leq k' \leq K} \ \big\Vert x_i - \mu_{k'} \big\Vert ,
\]
ties broken arbitrarily.

\emph{(Update.)} For a fixed $Z$, with associated groups $S_k = \{ i : z_{i,k} = 1\}$
assumed non-empty, the minimiser over $M$ is obtained independently for each cluster by
taking the mean of its points:
\[
\mu_k = \frac{1}{|S_k|} \sum_{i \in S_k} x_i
= \frac{\sum_{i=1}^n z_{i,k}\, x_i}{\sum_{i=1}^n z_{i,k}} .
\]
\end{lemma}

\begin{proof}
For the assignment step, write $\mathcal{L}(M, Z) = \sum_{i=1}^n \big( \sum_k z_{i,k}
\Vert x_i - \mu_k\Vert^2 \big)$. The constraint couples nothing across samples: each
row of $Z$ is chosen separately, and since that row must put its single $1$ somewhere,
the best place is the smallest of the $K$ numbers $\Vert x_i - \mu_k \Vert^2$.

For the update step, the objective splits as $\sum_{k=1}^K \big( \sum_{i \in S_k} \Vert
x_i - \mu_k \Vert^2 \big)$, and each bracket involves only $\mu_k$. So we must minimise
$\mu \mapsto \sum_{i \in S} \Vert x_i - \mu \Vert^2$ for a fixed set $S$. This is a
quadratic function of $\mu$, with Hessian $2 |S|\, I_d \succ 0$, hence strictly convex,
so its unique minimiser is the point where its gradient vanishes:
\[
\nabla_\mu \sum_{i \in S} \Vert x_i - \mu \Vert^2
= - 2 \sum_{i \in S} \big( x_i - \mu \big)
= 2 \Big( |S|\, \mu - \sum_{i \in S} x_i \Big) = 0
\quad \Longleftrightarrow \quad
\mu = \frac{1}{|S|} \sum_{i \in S} x_i ,
\]
which is the announced mean.
\end{proof}


\begin{center}
\fbox{\parbox{0.9\textwidth}{
\textbf{Lloyd's algorithm.} Initialise the centroids $\mu_1^{(0)}, \dots, \mu_K^{(0)}
\in \R^d$. Then for $t = 0, 1, 2, \dots$ repeat:
\begin{enumerate}
    \item \emph{Assignment.} $Z^{(t+1)} = \argmin_Z \ \mathcal{L}(M^{(t)}, Z)$: assign
    each sample to its nearest current centroid.
    \item \emph{Update.} $M^{(t+1)} = \argmin_M \ \mathcal{L}(M, Z^{(t+1)})$: replace
    each centroid by the mean of the samples just assigned to it.
\end{enumerate}
Stop when the assignment no longer changes.
}}
\end{center}

\paragraph{Why it stops.}
The convergence argument is unusual and pleasantly elementary: it is not an analytic
statement about small step sizes, but a counting argument.

\begin{prop}[Convergence in finitely many steps]
\label{prop:lloyd-converges}
The sequence $\mathcal{L}(M^{(t)}, Z^{(t)})$ is non-increasing and bounded below by $0$.
Moreover Lloyd's algorithm terminates after finitely many iterations.
\end{prop}

\begin{proof}
Each of the two steps minimises $\mathcal{L}$ in one variable with the other held fixed,
so neither can increase it:
\[
\mathcal{L}\big(M^{(t)}, Z^{(t)}\big) \ \geq \ \mathcal{L}\big(M^{(t)}, Z^{(t+1)}\big)
\ \geq \ \mathcal{L}\big(M^{(t+1)}, Z^{(t+1)}\big) ,
\]
and $\mathcal{L} \geq 0$ as a sum of squares. For termination, note that after the
update step the centroids are \emph{determined} by the assignment, so from $t \geq 1$
onwards the value $\mathcal{L}(M^{(t)}, Z^{(t)})$ is a function of $Z^{(t)}$ alone. There are at most
$K^n$ admissible assignment matrices, hence finitely many possible values. A
non-increasing sequence taking finitely many values is eventually constant; and as soon
as the value fails to decrease strictly, the assignment step has returned the same $Z$
and the algorithm has stopped by its own criterion.
\end{proof}

\begin{remark}[Stopping is not succeeding]
\Cref{prop:lloyd-converges} guarantees termination, and nothing more. The point reached
is a \emph{local} minimum, in the weak sense that neither step alone can improve it, and
it can be arbitrarily worse than the global one: it is easy to draw four obvious
clusters and an initialisation from which Lloyd's algorithm splits one of them in two
and merges another two, then stops, perfectly stable and perfectly wrong. This is
exactly the difficulty that convexity spared us throughout the supervised part of these
notes, and there is no way around it here.

The standard remedies are (i) to run the algorithm several times from different random
initialisations and keep the run with the smallest final $\mathcal{L}$, which is cheap
and is what libraries do by default, and (ii) to initialise cleverly rather than
uniformly at random, spreading the initial centroids out. The latter is the
\emph{$K$-means$++$} initialisation: pick the first centroid uniformly among the data
points, then each subsequent one at random with probability proportional to the squared
distance to the nearest centroid already chosen. It costs almost nothing and improves
results markedly.
\end{remark}

\paragraph{What can go wrong.}
Beyond the local minima, three limitations are worth knowing, all of them consequences
of the squared Euclidean distance in the objective.
\begin{itemize}
    \item \emph{The number of clusters must be given.} $K$ is a hyperparameter, and
    unlike in the supervised setting there is no held-out error to select it: the
    objective $\mathcal{L}$ decreases mechanically as $K$ grows, reaching $0$ at $K = n$,
    so it cannot be used to choose $K$ either. The usual device is again an
    \emph{elbow plot} of the final $\mathcal{L}$ against $K$, looking for the point
    where the curve flattens.
    \item \emph{The clusters come out round, and of comparable size.} Minimising squared
    distances to a centre means the implied shape is a ball, and the boundary between
    two clusters is the hyperplane bisecting the segment between their centroids. Two
    interleaved crescents, or one long thin cluster next to a round one, will be cut in
    a way no human would choose. If the natural groups are not blobs, $K$-means is the
    wrong tool.
    \item \emph{Outliers hurt twice.} A single far-away point contributes an enormous
    squared distance, so it drags the centroid of whichever cluster adopts it, and it may
    even claim a whole cluster for itself. Squared errors are not robust, exactly as in
    \Cref{sec:regression}.
\end{itemize}
To which one should add the familiar warning that the Euclidean distance mixes the
coordinates, so features on wildly different scales must be standardised first.

\begin{exercise}[$\star$]
Why does $\min \mathcal{L}$ decrease as $K$ increases? (Show that any solution with $K$
clusters can be turned into one with $K+1$ clusters and no larger loss.) Deduce that
choosing $K$ by minimising $\mathcal{L}$ would always return $K = n$.
\end{exercise}

\begin{exercise}[$\star\star$]
Take $n = 4$ points in $\R^2$ at the corners of a rectangle that is much wider than it
is tall, and $K = 2$. Find the global minimiser. Then exhibit an initialisation of the
centroids from which Lloyd's algorithm converges to the other, worse, partition.
\end{exercise}


